\chapter{Evaluating Adversarial Robustness: Benchmark Design and Metric Analysis}
\label{chap:advml}

% Revised under T-003 and D-041 from notes/chapter6-source-records/.
% D-030 fixes RQ4; D-031 fixes attribution; D-009 fixes the grey-box convention.

\section{Research Problem and Contribution}
\label{sec:advml:1}

This chapter addresses RQ4: \rqFourText{} The question concerns gap~\gapref{4},
which links adversary assumptions to the validity of the resulting measurements.

Adversarial machine learning studies how deliberate changes to training data or inference
inputs can cause a learning system to fail, together with methods for resisting such changes
\cite{biggio2018wildpatterns}. A poisoning attack intervenes during training, whereas an
evasion attack modifies an input presented to an already trained model. Defences may expose
the model to adversarial examples during training, transform its inputs or internal
representations, or detect and respond to suspected adversarial inputs. This chapter concerns
evasion attacks and the evaluation of these three defence strategies. The broader attack
taxonomy and terminology are defined in Section~\ref{sec:background:adversarial}.

An empirical robustness estimate depends on the permitted perturbations and the attacks used
to search for them \cite{biggio2014security, carlini2019evaluating}. An unsuccessful search
may reflect resistance to the tested attack, or a failure of the attack to exploit the model.
For example, Athalye et al. identified obstructed gradients in seven of nine non-certified
defences accepted at ICLR~2018 and circumvented six completely and one partially
\cite{athalye2018obfuscated}. Evaluators therefore need to examine the measurement procedure
alongside the defended classifier.

The SPARTA SAFAIR work produced two complementary evaluation mechanisms: an adversarial
machine-learning competition and a reusable benchmark framework
\cite{sparta2021d73,sparta2022d76}. The competition evaluated submitted defences on hidden
test samples. The benchmark provided common interfaces for evaluating trained models against
reference attack implementations, with task-specific components for different prediction
problems. Its versioned source code is publicly available on GitHub
\cite{norouzian2021adversarialbenchmark}.

Using the project records, this chapter reconstructs both evaluation mechanisms and assesses
whether the reported measurements are consistent with the documented attack objectives,
attacker access and perturbation constraints. The empirical evidence comprises the contest
submissions and the application of the benchmark to partner defences on histopathology images.
The chapter then derives properties of the accuracy-loss score and the binary detector response
policy. These analytical results establish what the published aggregates support and what
additional observations would be required to distinguish competing explanations. A short
comparison with the project's PDF-malware study examines how feature constraints alter the
evaluation problem.

SPARTA was funded under Horizon~2020 grant agreement No.~830892. Its seventh work package,
SAFAIR, concerned the security of artificial-intelligence systems. The candidate wrote the
whole of D7.6, which reports the contest and benchmark, and Chapter~5 of D7.5
\cite{sparta2022d76, sparta2021d75}. Section~\ref{sec:advml:10} distinguishes this contribution
from the partner methods and experiments discussed below.

\section{Threat Analysis and Evaluation Scope}
\label{sec:advml:2}

SPARTA's threat analysis links attack techniques to affected ML assets and the adversary's
knowledge \cite[Sections~3.3.2--3.3.6]{sparta2021d71}. Applied to the evaluations in this
chapter, these distinctions locate the attack within the ML life cycle. Poisoning concerns
training data or the learning process; evasion concerns inputs supplied to a trained model.

The evaluations examined here concern evasion. The contest specified attacks that induce a
chosen identity or errors in selected facial attributes \cite[Section~3.3]{sparta2021d73}.
The histopathology benchmark reports classification accuracy on clean and perturbed images
\cite[Section~3.1.3]{sparta2022d76}. The affected asset is therefore the inference input, and
the measured outcome is prediction correctness under attack. Sections~\ref{sec:advml:3}
and~\ref{sec:advml:4} specify attacker access and permitted perturbations, which determine
the conditions under which these measurements can be interpreted.

\section{Evaluation Requirements and Threat Models}
\label{sec:advml:3}

The threat model specifies the attack objective, the information available to the attacker and
the operations permitted on the input. Section~\ref{sec:background:adversarial} fixes the
knowledge terminology used throughout the dissertation. Under its grey-box convention, some
model information is known and the training process cannot be influenced. Knowledge of
training data and the ability to alter those data are separate assumptions.

The contest granted knowledge of the training corpus, withheld the defence technique and model
architecture, and allowed gradient-based updates against the deployed model
\cite{sparta2021d73, norouzian2021safaircontest}. Access to gradients is an oracle capability.
The term grey box alone therefore leaves the experimental interface underspecified. The
conditions stated in Section~\ref{sec:advml:4}, including the query interface and input budget,
define the contest adversary used here. This is an academic evaluation setting; a survey of
271 practitioners found that academic knowledge assumptions often grant more information than
attackers have in deployment \cite{grosse2024towards}.

D7.3 adopted the evaluation checklist of Carlini et al. It called for attacks adapted to the
complete defence, explicit attack hyperparameters and comparisons with transfer or
gradient-free attacks. It also requested perturbation-budget curves and ablations
\cite{sparta2021d73, carlini2019evaluating}. These requirements provide a basis for assessing
the delivered results. Participation by independent attack teams can broaden the search, but
the scoring rule and the validity of each generated input require separate checks.

\section{Tasks and Evaluation Protocol}
\label{sec:advml:4}

\subsection{Prediction Tasks and Permitted Perturbations}

The contest specified an attack and a defence track for each of two tasks
\cite{sparta2021d73}. In targeted face re-identification, the attacker receives an image and
a target identity and attempts to induce that identity as the prediction. In attribute
alteration, the attacker receives a nominated subset of the forty binary facial attributes
and attempts to cause errors on that subset. A participant could attack one task and defend
the other, but could not enter both sides of the same task.

Let $f_\theta$ denote the model scores and $h$ the resulting class decision. For an image $x$
in a specified input coordinate system, define the feasible perturbations by
\begin{equation}
  \mathcal S_\varepsilon(x)
  =\{x'\in\mathcal X_{\mathrm{valid}}:\|x'-x\|_\infty\leq\varepsilon\}.
  \label{eq:advml:feasible}
\end{equation}
Here $\mathcal X_{\mathrm{valid}}$ includes the permitted pixel range and quantisation.
For pixel values scaled to $[0,1]$, the contest budget was $\varepsilon=8/255$, corresponding
to a maximum change of eight intensity levels per channel in the discrete $0$--$255$ encoding.
The rules required discrete output images, deterministic and stateless submitted models, and
an upper limit of four hours for generation of 1,000 perturbed examples
\cite{sparta2022d76}. The report identifies two 12~GB NVIDIA GPUs, a Titan~X Pascal and a
GeForce~GTX, with a 32-core AMD CPU and 128~GB RAM. Container execution used Ubuntu~18.04
and CUDA~10.1. These are reported resource limits, not measured per-example runtimes.

The budget must be interpreted before any model-specific normalisation. For example, a linear
rescaling from $[0,1]$ to $[-1,1]$ doubles the numerical distance. A reproducible evaluation
therefore records both the input coordinate system and preprocessing. The contest's
$8/255$ setting is not assumed for the histopathology results, whose attack configurations
are absent from the reported result table.

\subsection{Scoring and the Delivered Protocol}

SAFAIR ranked submissions by the change between clean and perturbed accuracy,
\begin{equation}
  \Delta=A_{\mathrm{initial}}-A_{\mathrm{final}}.
  \label{eq:advml:delta}
\end{equation}
Attacks sought the largest $\Delta$ and defences the smallest. Unless stated otherwise,
accuracies and $\Delta$ are proportions in this chapter; $100\Delta$ is a difference in
percentage points. D7.5 specifies clean accuracy and runtime as tie-breakers
\cite{sparta2021d75}. These do not prevent a low-utility model with a smaller $\Delta$ from
outranking a more accurate model.

The planned opening round evaluated defences against FGSM, BIM and PGD, with respective
weights $0.2$, $0.4$ and $0.4$ \cite{sparta2021d73}.
FGSM makes one update in the sign of the input-loss gradient
\cite{goodfellow2015explaining}. BIM repeats smaller
FGSM-style updates and clips each result to the permitted input range and perturbation
neighbourhood \cite{kurakin2017bim}. PGD also applies iterative gradient updates, projects
each result onto the feasible perturbation set, and can use multiple starting points
\cite{madry2018pgd}. Carlini--Wagner instead
formulates adversarial-example generation as an optimisation problem that balances the
required classification outcome against perturbation size \cite{carlini2017robustness}.
The five selected attacks would then face the five selected defences, with the final score
averaged over opponents \cite{sparta2021d73}. No attack submissions arrived. The organisers
used fixed baseline attacks instead; D7.5 lists Carlini--Wagner in addition to FGSM, BIM and PGD
\cite{sparta2021d75}. Table~\ref{tab:advml:designdelivery} records the
change in protocol.

\begin{table}[tbp]
  \caption{Planned and reported SAFAIR evaluation procedures.}
  \label{tab:advml:designdelivery}
  \centering\small
  \begin{tabular}{@{}p{0.20\linewidth}p{0.36\linewidth}p{0.36\linewidth}@{}}
    \toprule
    & \textbf{Plan} \cite{sparta2021d73} & \textbf{Reported delivery} \cite{sparta2021d75,sparta2022d76} \\
    \midrule
    Participants & attack and defence tracks for both tasks & defence submissions only \\
    Attack selection & organiser baselines, then five selected participant attacks & four named organiser baselines \\
    Defence score & weighted opening $\Delta$; mean final $\Delta$ over opponents & aggregate best and worst perturbed accuracies; $\Delta$ from the best value \\
    Final comparison & five attacks against five defences & no participant-attack round \\
    Reproduction & hidden evaluation set and submitted containers & sampling code and tool descriptions; incomplete run-level records \\
    \bottomrule
  \end{tabular}
\end{table}

The lack of independent attacks reduced the diversity of the evaluation. It does not explain
all weaknesses of the resulting measurement. As derived in Section~\ref{sec:advml:8:metrics},
a constant classifier has $\Delta=0$ against any set of label-preserving attacks, including
the opponents in the planned final round. The scoring defect and the attack-coverage defect
are therefore distinct.

The contest drew inspiration from the NIPS~2017 competition, whose defence score was the
proportion of correctly classified attacked images averaged across attacks
\cite{kurakin2018competition}. That score differs from Equation~\ref{eq:advml:delta}.
The competitions also differed in task and participation, so their results do not isolate
the effect of independent contestants. Details of the schedule are retained in the cited
project records; the scientific comparison here concerns the scoring and attack interfaces.

\section{Data and Evaluation Measures}
\label{sec:advml:5}

\subsection{Corpora and Hidden Evaluation Samples}

The contest used CelebA for model training, a corpus described by Liu et al. as approximately
200,000 images of 10,000 identities, with forty binary attributes per image
\cite{liu2015celeba}. For hidden attribute evaluation the organisers used LFWA, the
forty-attribute annotation of Labeled Faces in the Wild, which contains 13,233 images of
5,749 identities \cite{huang2007lfw, liu2015celeba}. A separate corpus reduces direct image
reuse, but does not establish equality of image distributions or attribute prevalences.

Re-identification required identity labels compatible with the trained classifier. The
organisers therefore used CelebA images under a random choice of colour jitter, greyscale
conversion, horizontal reflection or rotation of at most five degrees
\cite{sparta2022d76}. The printed procedure applies a face-detection check and retries failed
transformations before excluding an image. The resulting selection may differ from an
unfiltered sample. The report recognises that participants could have used similar
augmentations during training. Augmentation alone does not demonstrate independence from
training images.

The sampling code in D7.6 calls for 1,000 LFWA examples and, separately, 1,000 examples from
a re-identification test CSV. These are two task-specific samples. The code assigns
\texttt{random.seed = 42} instead of calling a seeding function, while the two pandas samples
receive no explicit random state. The augmentation draws are also stochastic. Consequently,
the printed procedure does not regenerate the exact hidden images. Saved manifests and
transformed files could establish their identity, but their recovery is not documented in
the deliverable \cite{sparta2022d76}.

\begin{table}[tbp]
  \caption{Evaluation units and data provenance. Stated sample sizes describe the printed
    procedures, not recovered per-example result files.}
  \label{tab:advml:data}
  \centering\small
  \begin{tabular}{@{}p{0.23\linewidth}p{0.33\linewidth}p{0.36\linewidth}@{}}
    \toprule
    \textbf{Evaluation} & \textbf{Input and prediction} & \textbf{Source and sampling} \\
    \midrule
    Contest identity & one image; one identity label & augmented CelebA; 1,000 selected test records \cite{sparta2022d76} \\
    Contest attributes & one image; forty binary labels & LFWA; 1,000 selected images \cite{sparta2022d76} \\
    Partner benchmark & one $50\times50$ histopathology patch; IDC/non-IDC & training, validation and test subsets described; evaluated subset size and patient separation unreported \cite{sparta2020d72} \\
    PDF comparison & 21 tag counts; benign/malicious & 20,000 files; 45/45/10 per cent defender/proxy/crafting split \cite{sparta2021d75} \\
    \bottomrule
  \end{tabular}
\end{table}

The histopathology corpus is identified in D7.2 as the IDC patch collection, whereas D7.6
calls it a breast-cancer image dataset without naming it \cite{sparta2020d72,sparta2022d76}.
The full collection contains 198,738 negative and 78,786 positive patches, so its majority
share is approximately $0.7161$ \cite{mooney2017idc}. The test subset's class share is not
reported. Neither this collection-wide ratio nor the contest's sample sizes can be used to
reconstruct the benchmark's confusion matrices.

\subsection{Accuracy and Task-Specific Attack Success}

For $N$ examples with identity labels $y_i$, the decision is
$h(x_i)=\arg\max_k f_{\theta,k}(x_i)$. Attribute decisions are
$h_j(x_i)=\mathbf1[f_{\theta,j}(x_i)\geq0]$ when the outputs are logits, with equivalent
threshold $0.5$ for sigmoid probabilities. The corresponding accuracies are
\begin{equation}
  A_{\mathrm{id}}=\frac1N\sum_{i=1}^N\mathbf1[h(x_i)=y_i],
  \qquad
  A_{\mathrm{attr}}=\frac1{40N}\sum_{i=1}^N\sum_{j=1}^{40}
                   \mathbf1[h_j(x_i)=y_{ij}].
  \label{eq:advml:accuracies}
\end{equation}
These conventions follow the task-specific scoring code printed in D7.6. Attribute accuracy
counts individual attribute decisions; it does not require all forty predictions for an
image to be correct \cite{sparta2022d76}.

Accuracy loss does not identify success against a specified target. For identity targets
$t_i\neq y_i$, let $I_t=\{i:h(x_i)=y_i,\ t_i\neq y_i\}$ contain eligible, initially correct
examples. Given valid attacked inputs $x_i'$, conditional targeted success is
\begin{equation}
  \mathrm{ASR}_{\mathrm{target}}
  =\frac{1}{|I_t|}\sum_{i\in I_t}\mathbf1[h(x_i')=t_i].
  \label{eq:advml:targeted}
\end{equation}
A prediction of a third identity reduces accuracy without satisfying the targeted objective.
An empty eligible set makes this rate undefined and its size must be reported.

For a nonempty nominated attribute subset $J_i$, define
$I_J=\{i:h_j(x_i)=y_{ij}\text{ for all }j\in J_i\}$. One explicit success criterion is
\begin{equation}
  \mathrm{ASR}_{\mathrm{subset}}
  =\frac1{|I_J|}\sum_{i\in I_J}\prod_{j\in J_i}
    \mathbf1[h_j(x_i')\neq y_{ij}].
  \label{eq:advml:subset}
\end{equation}
This requires errors on every nominated attribute of an initially correct subset. Reporting
$|I_J|$ separates induced errors from those already present; the rate is undefined if
$|I_J|=0$. Equations~\ref{eq:advml:targeted}
and~\ref{eq:advml:subset} specify measurements needed to evaluate the stated objectives;
they are not additional scores recovered from the contest. If an attack flips all $k$ initially
correct bits in the nominated subset and leaves all other decisions unchanged, the per-image contribution
to attribute accuracy falls by only $k/40$. For $k=1$ this is 2.5 percentage points, although
the specified attack succeeds.

\subsection{Aggregation over Attacks}

An attack suite can fail on different examples under each method. Selecting the lowest
aggregate accuracy does not account for that complementarity. Following per-example
evaluation guidance \cite{carlini2019evaluating}, let $\mathcal A_\varepsilon$ be a finite
set of attacks whose outputs satisfy Equation~\ref{eq:advml:feasible}. For a single-label task,
define
\begin{equation}
  \widehat A_{\mathcal A}(\varepsilon)
  =\frac1N\sum_{i=1}^N
   \min_{a\in\mathcal A_\varepsilon\cup\{\mathrm{id}\}}
   \mathbf1[h(a(x_i))=y_i].
  \label{eq:advml:suite}
\end{equation}
The identity transformation includes clean errors in the score. The suite must respect the
same threat model for each defence. Because it searches only a subset of valid perturbations,
its surviving accuracy is an upper bound on worst-case accuracy on these examples, not a
certificate of robustness. Moreover,
\begin{equation}
  \widehat A_{\mathcal A}(\varepsilon)
  \leq\min_{a\in\mathcal A_\varepsilon\cup\{\mathrm{id}\}}
           \frac1N\sum_i\mathbf1[h(a(x_i))=y_i],
  \label{eq:advml:aggregation-bound}
\end{equation}
since each per-example minimum is no greater than its value for any fixed attack.
The aggregate tables do not identify which examples fail under which attack, so this suite
score cannot be calculated retrospectively from them.

\subsection{Biometric Data and Data Protection}
\label{sec:advml:5:ethics}

The contest processed photographs of identifiable people and predicted identities or facial
attributes. D7.1 discusses purpose limitation, data minimisation and privacy by design, but
the contest documentation supplies no corpus-specific account of the lawful basis or the
participants' data-governance responsibilities \cite{sparta2021d71,sparta2022d76}. The
absence of that documentation establishes a reporting gap; it does not establish the consent
status of each person depicted or a legal finding about every possible use.

Attribute imbalance also affects the scientific interpretation. The SD account describes a
constant attribute predictor whose accuracy remains high because some attributes are rare
\cite{sparta2022d76}. A complete evaluation would therefore report attribute prevalences and
per-attribute outcomes alongside the aggregate score. Demographic performance differences
cannot be inferred from the reported aggregates. This chapter confines its empirical claims
to the documented classification tasks and does not infer suitability for biometric deployment.

\section{Benchmark Architecture and Reproducibility}
\label{sec:advml:6}

The benchmark implements a common software procedure for adversarial evaluation. It separates
attack execution from the task-specific model and data interfaces. D7.3 proposed five modules,
which D7.6 reports as implemented: attacks, datasets and models, together with use cases and a
wrapper \cite{sparta2021d73,sparta2022d76}. Figure~\ref{fig:advml:architecture} redraws the
published architecture.

\begin{figure}[tbp]
  \centering
  \includegraphics{sparta/thesis/sparta-benchmark-architecture.pdf}
  \caption{Benchmark architecture, redrawn after D7.6 Figure~2, p.~22 \cite{sparta2022d76}.
    Box labels follow the original. The second line in each box paraphrases the module and
    run descriptions in D7.6 Section~3.1.1. The five modules separate task definitions and
    data loading from attack execution. The diagram records the software design; it does not
    establish that a particular experiment preserved its model weights, data manifest and
    attack settings.}
  \label{fig:advml:architecture}
\end{figure}

Within this framework, standardisation concerns the component interfaces and execution
procedure. A run loads a trained checkpoint and evaluation data, converts the model to the
attack-library interface, measures performance before and after the selected attacks, and emits
the accuracies in one JSON-formatted result. Researchers can extend the procedure through new
attack classes and task-specific components. The implementation of this reusable workflow is
available in the versioned GitHub repository
\cite{norouzian2021adversarialbenchmark}.

The wrapper exposes a model through Foolbox, whose reference implementations supply the
attacks \cite{rauber2017foolbox}. PyTorch is the native model interface; the report also
describes conversion of TensorFlow or Keras models. Use-case classes select the prediction
criterion and loss for the task. A new attack extends a base class and is discovered by the
loader; a configuration list excludes named attacks. These interfaces permit reuse across
models. Compatibility still depends on the output semantics: the report's multi-label
wrapper handles fixed-budget attacks, while some minimum-distance interfaces are omitted
\cite{sparta2022d76}.

D7.6 also contains a description of splitting data and training classifiers, alongside repeated
statements that weights must already exist. The documented invocations require checkpoints,
and the inspected public implementation loads trained models
\cite{norouzian2021adversarialbenchmark}. Accordingly, the benchmark is described here as an
evaluation tool. Its public source is separate from the contest development toolkit
\cite{norouzian2021safairtoolkit}.

\paragraph{Perturbation validation.} The public revision identified by commit
\texttt{2fac62b947ed} pins Foolbox~3.3.1. Its attack call returns a raw candidate and a
budget-clipped candidate, with a success indicator. However, lines~134--135 of
\texttt{main.py} pass the raw candidate to the accuracy function
\cite{norouzian2021adversarialbenchmark}. In Foolbox's minimisation interface the clipped
candidate is obtained by restricting the raw perturbation to the requested distance budget;
the raw output need not satisfy that budget \cite{foolbox2021api}. Scoring the raw output
can therefore count changes outside the permitted set. This is a finding about the inspected
implementation. No recovered run record connects that revision to the histopathology table,
so it does not establish the cause of any historical result.

A further distinction concerns targeted runs. The same implementation replaces the reference
labels with the nominated target before computing its fields labelled accuracy. Those fields
then measure target-label agreement, with a different interpretation from true-label accuracy
\cite{norouzian2021adversarialbenchmark}. The examples in D7.6 Figure~6 likewise show a
targeted result whose value rises under attack \cite{sparta2022d76}. Each output therefore
needs an explicit task, reference label and denominator.

The inspected source fixes several dependency versions and allows the implementation to be
identified by a commit. Reproduction of a measured result additionally requires the exact
checkpoint and evaluation manifest, the preprocessing and attack parameters, and the individual
outputs. A valid evaluation would check quantisation and distances after all conversions,
record failed or skipped attacks, and score the candidate that satisfies the declared budget.
These are requirements derived from the interfaces, not claims that the historical runs
performed those checks.

\section{Defence Mechanisms and Integration}
\label{sec:advml:7}

The evaluated configurations instantiate the three defence strategies introduced in
Section~\ref{sec:advml:1}: adversarial training, transformations of inputs or internal
representations, and adversarial-input detection.

D7.6 evaluates five partner configurations described in D7.2. The underlying histopathology
classifier combines a pretrained VGG16 feature extractor with a dense prediction network
\cite{sparta2020d72,sparta2022d76}. Write its scores as
$f_\theta(x)=c_\psi(\phi_\omega(x))$, where $\phi_\omega$ is the convolutional feature
extractor and $c_\psi$ the prediction head. The following notation makes the intervention
of each defence explicit without supplying missing layer configurations.

\subsection{Training and Representation Transformations}

The adversarial-training configuration retrains the classifier with clean and perturbed
examples. D7.6 identifies FGSM, BIM and PGD as the attacks used to generate its adversarial
training examples \cite{sparta2022d76}. The surviving description does not specify the
mixture ratio or a complete optimisation schedule. This method exposes the model to those
attacks during training; resistance to other attacks remains an empirical question.

For the dimensionality-reduction variants, let $r_\eta$ denote an autoencoder. The two
configurations have score functions
\begin{equation}
  f_{\mathrm{top}}(x)=c_\psi\bigl(\phi_\omega(r_{\eta_x}(x))\bigr),
  \qquad
  f_{\mathrm{mid}}(x)=c_\psi\bigl(r_{\eta_z}(\phi_\omega(x))\bigr).
  \label{eq:advml:compositions}
\end{equation}
The first transforms pixels and the second transforms extracted features.
D7.2 states that the original convolutional network and prediction head retain their weights
in these autoencoder variants. Its separate encoder-only variant trains a new prediction
head, but that variant is not a column of the D7.6 benchmark table
\cite{sparta2020d72}. Figure~\ref{fig:advml:autoencoders} redraws the two published
placement diagrams in this notation.

\begin{figure}[tbp]
  \centering
  \includegraphics{sparta/thesis/sparta-autoencoder-placement.pdf}
  \caption{Autoencoder placement in the partner classifier, redrawn after D7.2 Figures~13
    and~11, p.~14 \cite{sparta2020d72}. In the top variant (a), the autoencoder
    $r_{\eta_x}$ transforms the input before feature extraction. In the middle variant (b),
    $r_{\eta_z}$ acts between the feature extractor $\phi_\omega$ and the prediction head
    $c_\psi$, as in Equation~\ref{eq:advml:compositions}. The legend follows D7.2, which
    keeps the original weights of VGG16 and the dense network in both variants. The
    diagrams specify placement, not a matched ablation of all training choices.}
  \label{fig:advml:autoencoders}
\end{figure}

A smaller reconstruction error alone does not establish robustness of either composition.
The relevant attack must optimise against the complete transformed classifier, including its
preprocessing. Reusing examples generated against the original classifier tests a different
condition from attacking the defended composition \cite{carlini2019evaluating}.

\subsection{History and Activation Detectors}

Figure~\ref{fig:advml:detectors} shows the two partner detectors and their response paths.

Prediction similarity adds a history of queries and their predictions. The partner study
records user information with image similarity and alarm statistics, and reports SSIM as the
similarity measure \cite{sparta2020d72}. Its alarm can be written
$d_t=d(x_t,H_{t-1})$, where $H_{t-1}$ contains the preceding queries associated with the
monitored interaction. A sequence of related optimisation queries may then trigger an alarm.
A transferred adversarial example submitted once presents a different observation to this
stateful detector.

For reproducibility, a history-based evaluation must specify the sequence order, user
association and reset policy. The contest's stateless-submission rule does not describe this
benchmark configuration. The reports do not establish a common history protocol across the
attacks in the result table, so their differences cannot be attributed solely to whether an
attack uses gradients.

The activation detector instead characterises the response of the underlying classifier to
an input. For each of nineteen hidden layers it calculates nine statistics: mean and median;
skewness and kurtosis; standard deviation and interquartile range; minimum and maximum;
and the number of nonzero activations \cite{sparta2020d72}. With layer activations
$a_\ell(x)$ and the statistic vector $s$, the representation is
\begin{equation}
  z(x)=\operatorname{concat}_{\ell=1}^{19}s(a_\ell(x))\in\mathbb R^{171}.
  \label{eq:advml:activation-features}
\end{equation}
After standardisation and feature selection, two support-vector classifiers are trained, one
for each class predicted by the original model. The selected detector therefore depends on
that prediction. The 94\% detector accuracy reported in D7.2 concerns detection of adversarial
inputs; it is not the defended task accuracy of the D7.6 table.

\begin{figure}[tbp]
  \centering
  \includegraphics{sparta/thesis/sparta-detector-mechanisms.pdf}
  \caption{Partner detection mechanisms, redrawn after D7.2 Figures~15 and~16, p.~16
    \cite{sparta2020d72}. Panel (a) uses preceding queries, whereas panel (b) uses internal
    activations of the classifier. The recorded parameters, the activation features of
    Equation~\ref{eq:advml:activation-features} and the per-class detector follow D7.2.
    The response is the label reversal of Equation~\ref{eq:advml:reversal}, which D7.6
    applies in the benchmark \cite{sparta2022d76}. The final action is part of the defended
    system and must be included when task accuracy is measured.}
  \label{fig:advml:detectors}
\end{figure}

\subsection{From an Alarm to a Classification}

To place the detectors in an accuracy table, D7.6 uses the binary response policy of returning
the opposite class whenever the detector raises an alarm \cite{sparta2022d76}. Let
$h(x)\in\{0,1\}$ denote the base prediction, $d(x)\in\{0,1\}$ the alarm and $g(x)$ the
returned class. For a fixed history, where applicable,
\begin{equation}
  g(x)=h(x)\oplus d(x),
  \label{eq:advml:reversal}
\end{equation}
where $\oplus$ denotes exclusive OR. This rule corrects an erroneous binary prediction when
the detector raises an alarm, but turns a correct prediction into an error under the same
alarm. Consequently, on a common labelled evaluation set,
\begin{equation}
  A(g)-A(h)
  =\Pr(d=1,h\neq y)-\Pr(d=1,h=y).
  \label{eq:advml:reversal-gain}
\end{equation}
The probabilities denote empirical fractions on that set. To obtain the identity, partition
the observations by the alarm and the correctness of $h$: unflagged predictions are unchanged,
flagged errors are corrected, and flagged correct predictions are reversed. If $d(x)=1$ for
every observation, then $A(g)=1-A(h)$, while $g$ can still depend on the input.

Detection accuracy alone does not determine either term in
Equation~\ref{eq:advml:reversal-gain}. A useful evaluation records alarms jointly with base
and returned predictions, and measures false alarms on clean inputs. Abstention would be a
different response policy, for which coverage and error among accepted predictions must be
reported \cite{carlini2019evaluating}. The historical benchmark reports the reversal policy;
its aggregate accuracies do not identify the two joint fractions.

\section{Results and Metric Analysis}
\label{sec:advml:8}

\subsection{Contest Evidence}
\label{sec:advml:8:contest}

Table~\ref{tab:advml:contest} preserves the numerical columns of D7.6 Table~1. The source
reports one row per ranked team and task, rather than one row per submitted checkpoint.
Its surrounding text describes best and worst scores across submissions, but does not link
each value to an attack and checkpoint \cite{sparta2022d76}.

\begin{table}[tbp]
  \caption{Reported contest scores, including the worst-perturbed column, from D7.6 Table~1,
    p.~16 \cite{sparta2022d76}. Task names follow that table. The SD task label conflicts
    with its detailed attribute-model description; best and worst values are not identified
    as paired outcomes for one checkpoint under a common attack suite.}
  \label{tab:advml:contest}
  \centering\small
  \begin{tabular}{@{}llp{0.24\linewidth}cccc@{}}
    \toprule
    Rank & Team & Task & Clean & Best & Worst & $\Delta$ \\
    \midrule
    1 & SD & re-identification & 0.87 & 0.85 & 0.67 & 0.02 \\
    2 & Vicomtech & re-identification & 0.85 & 0.72 & 0.72 & 0.13 \\
      & Vicomtech & attribute alteration & 0.80 & 0.65 & 0.65 & 0.15 \\
    3 & BPI & re-identification & 0.84 & 0.68 & 0.68 & 0.16 \\
    \bottomrule
  \end{tabular}
\end{table}

The SD row reports 0.85 as best perturbed accuracy and 0.67 as worst, a spread of eighteen
percentage points. The printed $\Delta=0.02$ uses the best value. Subtracting 0.67 from 0.87
would give 0.20, but it would not establish a worst-case loss for the same model without the
missing checkpoint association. The report's prose also substitutes 87\% for attacked
accuracy and gives 61\% and 88\% values that do not match its table. The numerical record
used here is the table, with those inconsistencies retained as limits on interpretation.

The counts of submissions also disagree. D7.5 states seven defence submissions, six for
re-identification and one for attributes, whereas D7.6 states six submissions in total
\cite{sparta2021d75,sparta2022d76}. An unscored ITTI account appears in D7.6, but the
reports do not state that this explains the difference. Similarly, D7.5 gives SD a zero
$\Delta$ while D7.6 Table~1 gives 0.02. The development accounts describe models compatible
with different scores, but do not supply the identifiers needed to reconcile the two reports.

The four rows do not establish a general relation between clean performance and robustness.
They combine different prediction tasks and selected submissions from three teams. No
regression or significance test is appropriate for this table. D7.6 mentions ten-fold
cross-validation in the Vicomtech and BPI accounts, but does not provide fold assignments or
fold-level outcomes. Those statements cannot supply uncertainty estimates for the printed
scores.

\subsection{SD Development Configurations}
\label{sec:advml:8:sd}

D7.6 describes three SD models. Its introductory text assigns them to re-identification,
while the model details specify forty attribute outputs. The high accuracy of a constant
attribute vector cannot be interpreted as a demonstrated identity-classification result.
The following analysis therefore follows the attribute-model description and leaves its
mapping to Table~\ref{tab:advml:contest} unresolved \cite{sparta2022d76}.

The first configuration was a convolutional network trained on PGD examples with
$\varepsilon=0.3$, approximately 9.6 times the contest's $8/255$ budget if both use
$[0,1]$ pixels. The team reports that training produced the same output for every input,
with approximately 80\% clean attribute accuracy and zero accuracy loss. It attributes that
accuracy to the prevalence of absent attributes and recognises a weakness in the dataset
and contest rules. This is the source's direct account of prediction collapse, although
neither the checkpoint nor individual predictions accompany the reported result.

The second configuration combines six EfficientNet and three pruned ECA-ResNet models.
Their averaged outputs are thresholded to obscure gradients. For a nonzero scalar attribute
logit $u$, write the transformation as
\begin{equation}
  q(u)=\begin{cases}10,&u>0,\\-10,&u<0.\end{cases}
  \qquad
  \mathbf1[q(u)\geq0]=\mathbf1[u\geq0],
  \qquad q'(u)=0\ \ (u\neq0).
  \label{eq:advml:threshold}
\end{equation}
The source does not specify the exact-zero tie. Away from that tie, the operation preserves
the binary decision and suppresses its direct derivative. It can obstruct a gradient-based
attack without changing the classifier's decision regions. Attacks adapted to the
non-differentiable operation or transferred from another model remain relevant
\cite{athalye2018obfuscated, tramer2020adaptive}.

The second configuration achieved approximately 92\% clean accuracy and 79\% under
transferred attacks in the team's development account. The third used one pruned ECA-ResNet,
added adversarial training, and retained the threshold operation. Its reported clean
accuracy was approximately 87\%, with 84--85\% under transferred attacks
\cite{sparta2022d76}. Figure~\ref{fig:advml:sd} shows these values separately from the
contest table.

\begin{figure}[tbp]
  \centering
  \includegraphics[width=\linewidth]{sparta/thesis/sparta-sd-comparison.pdf}
  \caption{SD development configurations from D7.6 Sections~2.8.1.2--2.8.1.3, pp.~18--19
    \cite{sparta2022d76}. The interval 84--85\% is a reported approximate range, not a
    confidence interval. The source describes attribute outputs despite its re-identification
    heading. These values are not additional official contest rows.}
  \label{fig:advml:sd}
\end{figure}

The third configuration trades five percentage points of reported clean accuracy for five to
six points of transferred-attack accuracy relative to the second. This is a descriptive
comparison between systems: ensemble size and training changed together, and the attacks'
individual settings are not given. It does not isolate the effect of adversarial training.
The transfer results nevertheless show why resistance to a direct gradient attack is an
incomplete measure of these constructions.

\subsection{Histopathology Benchmark}
\label{sec:advml:8:benchmark}

D7.6 reports that Vicomtech used the benchmark tool to evaluate five defence configurations
on a histopathology classification task \cite[Section~3.1.3]{sparta2022d76}. This application
documents use of the framework for a prediction problem distinct from the competition's facial
image tasks. Table~\ref{tab:advml:benchmark} reproduces the reported accuracies for binary IDC
classification. For the two detectors, these include the response policy in
Equation~\ref{eq:advml:reversal}. The clean row records the prediction cost associated with
each configuration.

\begin{table}[tbp]
  \caption{Published accuracies for the five partner defence configurations, from D7.6
    Table~2, pp.~28--29 \cite{sparta2022d76}. Attack hyperparameters and the evaluated
    subset are unspecified. Each entry is a reported aggregate, with no dispersion estimate.}
  \label{tab:advml:benchmark}
  \centering\small
  \begin{tabular}{@{}lccccc@{}}
    \toprule
    Attack & Adv.\ training & Top AE & Middle AE & Pred.\ similarity & Activations \\
    \midrule
    none (clean) & 0.844 & 0.726 & 0.817 & 0.844 & 0.708 \\
    FGSM & 0.837 & 0.716 & 0.812 & 0.837 & 0.708 \\
    $L_2$ basic iterative & 0.844 & 0.725 & 0.817 & 0.842 & 0.708 \\
    BIM & 0.837 & 0.716 & 0.812 & 0.837 & 0.708 \\
    DeepFool & 0.003 & 0 & 0.001 & 0.015 & 0.708 \\
    additive noise & 0.844 & 0.725 & 0.817 & 0.844 & 0.708 \\
    NewtonFool & 0.464 & 0 & 0 & 0.460 & 0.708 \\
    PGD & 0.837 & 0.717 & 0.813 & 0.837 & 0.708 \\
    \bottomrule
  \end{tabular}
\end{table}

The middle autoencoder exceeds the top autoencoder by 9.1 percentage points on clean data
and 9.6 points under the reported PGD condition. Its PGD loss is 0.4 points, compared with
0.9 for the top variant. These values favour the middle configuration within the reported
evaluation. They do not establish that placement alone caused the difference, because the
complete training configurations and repeated-run outcomes are unavailable.

Adversarial training and prediction similarity have equal reported clean accuracy and equal
FGSM, BIM and PGD values. Prediction similarity is higher under DeepFool by 1.2 percentage
points, but lower under NewtonFool by 0.4 points. It therefore does not dominate adversarial
training across the table. The source's explanation that similarity detects gradient-based
construction is insufficient to explain the results, since DeepFool and NewtonFool also use
gradient information. Query sequences and detector alarms would be needed to test that
mechanism.

The activation column remains at 0.708, including on clean inputs. This establishes equality
of the printed aggregates to their reported precision. It does not establish identical
predictions, an always-active detector or insensitivity to inputs. The full corpus's majority
share is approximately 0.7161, but the unknown evaluation split prevents a matched baseline
comparison. Equation~\ref{eq:advml:reversal-gain} also shows that always reversing predictions
would not, by itself, yield the class prior. The mechanism behind the column remains
unidentified.

The near-zero DeepFool results for four configurations require a different qualification.
Minimum-distortion attacks search for an adversarial input and should be compared by the
distortion required, or by success after imposing a common valid budget
\cite{moosavi2016deepfool,carlini2019evaluating}. Neither the norm and stopping settings nor
achieved distortions are reported for these benchmark runs. The table therefore cannot
establish that its attack rows represent equally constrained adversaries. The code finding
in Section~\ref{sec:advml:6} reinforces the need to validate distances, but does not resolve
the historical configuration. The additive-noise row is a useful control condition; without
its distribution, budget and number of trials, its small changes do not establish a
comprehensive random-search check.

\subsection{Sensitivity to the Scoring Rule}
\label{sec:advml:8:metrics}

Figure~\ref{fig:advml:deltas} subtracts each attacked accuracy from that configuration's
clean accuracy. It makes the score dependence visible without changing the original results.
Under PGD, the activation detector has the smallest $\Delta$, zero, but its attacked
accuracy is 12.9 percentage points below adversarial training. Similarly, the middle
autoencoder has a smaller loss than adversarial training while its attacked accuracy is
2.4 points lower. A smaller loss alone therefore does not identify the more useful defended
classifier.

\begin{figure}[tbp]
  \centering
  \includegraphics[width=\linewidth]{sparta/thesis/sparta-benchmark-deltas.pdf}
  \caption{Accuracy losses computed from D7.6 Table~2 \cite{sparta2022d76}. Each cell is
    $100(A_{\mathrm{clean}}-A_{\mathrm{attack}})$ in percentage points; the original
    accuracies remain in Table~\ref{tab:advml:benchmark}. The separation of attack rows
    does not imply matched budgets. In particular, zero loss for the activation detector
    describes its aggregate score, not a demonstrated mechanism or robustness certificate.}
  \label{fig:advml:deltas}
\end{figure}

For paired clean and attacked examples, let $C$ denote an initially correct prediction and
$C'$ a correct prediction after attack. Expanding the two accuracy terms yields
\begin{equation}
  \Delta=\Pr(C\cap\neg C')-\Pr(\neg C\cap C').
  \label{eq:advml:cancellation}
\end{equation}
Both accuracies contain the correctly classified observations in $C\cap C'$, which cancel.
Thus, $\Delta$ measures the net change in correctness. It is not the conditional fraction
of initially correct examples made incorrect. Equal numbers of harmful and corrective
changes can leave $\Delta=0$ despite changes in individual outcomes. The aggregate table
cannot determine those two counts.

\paragraph{Constant-prediction result.} For any constant classifier $h_c(x)=c$ and
label-preserving attack $a$, $h_c(a(x_i))=h_c(x_i)$ for each observation. Therefore,
\begin{equation}
  A_{\mathrm{clean}}(h_c)=A_{\mathrm{attack}}(h_c)
  =\widehat p_c,\qquad \Delta(h_c)=0,
  \label{eq:advml:constant}
\end{equation}
where $\widehat p_c$ is the evaluated class fraction. The statement holds for each attack
and any weighted average over attacks. Independent opponents cannot eliminate this property
of the score. The result establishes an indifference to prediction utility when ranking
only by loss; it does not imply that the constant classifier correctly classifies the other
classes.

For attributes, choose a fixed vector $m$ using training-set majority labels. Its held-out
accuracy is
\begin{equation}
  A_{\mathrm{attr}}(m)
  =\frac1{40}\sum_{j=1}^{40}
    \bigl[m_j\widehat p_j+(1-m_j)(1-\widehat p_j)\bigr],
  \qquad \Delta(m)=0,
  \label{eq:advml:constant-attributes}
\end{equation}
where $\widehat p_j$ is the prevalence of attribute $j$ in the evaluation set. When the
training and evaluation majorities agree, each summand is
$\max(\widehat p_j,1-\widehat p_j)$. This explains how a fixed attribute vector can have
high average accuracy on an imbalanced corpus. The baseline labels must be fixed before
evaluation; selecting them from hidden test labels would use unavailable information.

These identities support reporting clean utility together with attacked performance and a
matched constant baseline. Prediction distributions and paired error transitions provide
additional checks. None of these quantities can be reconstructed uniquely from an unchanged
accuracy column, so the activation result remains an observation rather than a second
demonstration of the SD collapse mechanism.

\subsection{Comparison with Constrained PDF-Malware Features}
\label{sec:advml:8:pdf}

The public D7.5 report contains a partner study of evasion against PDF-malware classifiers
\cite{sparta2021d75}. It provides a comparison with the continuous image setting, rather
than an additional experiment on the Android systems in this dissertation. The study uses
10,000 benign and 10,000 malicious Contagio files. PDFiD extracts 21 tag counts; the data
are partitioned into 45\% for defender training, 45\% for proxy training and 10\% for
adversarial crafting. A neural proxy generates examples transferred to neural, SVM and
random-forest classifiers, without queries to the target during generation.

The allowed changes increase tag counts and preserve integer values after inverse
standardisation and rounding. If $v\in\mathbb N_0^{21}$ is the original count vector,
the reported feature constraints require
\begin{equation}
  v'\in\mathbb N_0^{21},\qquad v'_j\geq v_j\quad(j=1,\ldots,21).
  \label{eq:advml:pdf-constraints}
\end{equation}
A perturbation bound also applies in the attack representation. Its interpretation depends
on standardisation and the subsequent rounding. These constraints restrict the feature
search, but do not establish that a corresponding modified PDF was constructed or that its
malicious functionality was preserved. The distinction is the feature-space versus
problem-space distinction of Section~\ref{sec:background:adversarial}.

\begin{figure}[tbp]
  \centering
  \begin{subfigure}[t]{0.49\linewidth}
    \centering
    \includegraphics[width=\linewidth]{sparta/report-originals/d75-fig15-iter-fgsm-left.jpg}
    \caption{Adversarial training.}
  \end{subfigure}\hfill
  \begin{subfigure}[t]{0.49\linewidth}
    \centering
    \includegraphics[width=\linewidth]{sparta/report-originals/d75-fig15-iter-fgsm-right.png}
    \caption{Feature scattering with adversarial training.}
  \end{subfigure}
  \caption{Original iterative-FGSM transfer curves from the partner PDF-malware study,
    D7.5 Figure~15, p.~23 \cite{sparta2021d75}. Both panels retain the published legends
    and values. The horizontal axis is the source's $\log_{10}(\delta)$ perturbation
    budget, distinct from the accuracy loss $\Delta$ used in this chapter. These are
    feature-space results, with no uncertainty bands or demonstrated construction of
    modified PDFs with preserved malicious functionality.}
  \label{fig:advml:pdf-transfer}
\end{figure}

The retained legend names the SVM, random forest (FST), neural network (ANN) and their
ensemble (Ensbl); the prefix r-fgsm denotes adversarial retraining with FGSM examples.
Proxy measures success against the attack's source model, rather than cross-model transfer.
In the feature-scattering condition the SVM and random forest receive adversarial
augmentation, so that panel represents a hybrid defence comparison \cite{sparta2021d75}.

In Figure~\ref{fig:advml:pdf-transfer}, transfer success for many reported models rises
sharply as the budget increases, and several approach 100\% at larger budgets. Other curves
are non-monotone. Such curves record the success of the specified search at each setting;
they do not establish the monotone worst-case envelope over nested feasible sets. The two
panels also show dependence on the target model and defence configuration. They support
reporting a budget-dependent curve with explicit constraints, while the limited settings
and absent repeated-run information prevent a general ranking of defence families. The
same methodological requirement applies to Android: an evaluation must state how permitted
artefact transformations map to the representation used by the detector.

\section{Threats to Validity}
\label{sec:advml:9}

\paragraph{Measurement validity.} The contest combines task-dependent accuracies with a
net-change score. The identity and attribute objectives require different success criteria,
and the targeted software output uses target-label agreement. Aggregate accuracy loss
cannot identify induced errors, useful classification or detector alarms. The analytical
results in Sections~\ref{sec:advml:7} and~\ref{sec:advml:8:metrics} identify these limits;
they do not recover unreported individual outcomes.

\paragraph{Attack coverage and budgets.} No independent attack submissions entered the
contest. Its fixed gradient-based baselines did not constitute a documented adaptive
assessment of every defence. The SD development account includes transfer attacks, but
the official table provides no per-attack breakdown. The benchmark additionally mixes
attack families without the full configurations needed for comparison. A later evaluation
would require valid perturbations in common coordinates and attacks adapted to the complete
defended system, with convergence checks \cite{carlini2019evaluating,tramer2020adaptive}.

\paragraph{Sampling and uncertainty.} The printed hidden-set procedure has unseeded draws
and stochastic transformations. Augmentations do not establish absence of overlap with
training images. The histopathology subset size, class proportions and patient-level
partitioning are unreported. Patch counts cannot be treated as independent patient counts,
and forty attributes on the same image are dependent observations. The available rounded
aggregates and statements about cross-validation do not support confidence intervals,
paired significance tests or a fitted clean--robustness relation.

\paragraph{Conflicting records and implementation provenance.} The SD task names and the
submission totals disagree across the reports. No checkpoint mapping resolves the zero
and 0.02 deltas or the best and worst score associations. The inspected code exposes a
raw-versus-clipped scoring issue, but its connection to the published experiment is unknown.
The source repository and dependency versions therefore support implementation inspection,
while experiment reproduction still requires the corresponding data and model artefacts.
These distinctions remain explicit rather than being resolved by a presumed run history.

\paragraph{Generalisation and attribution.} The four participating teams belonged to SPARTA,
so the contest did not test participation beyond the consortium. The evaluated image tasks
also differ from Android malware analysis in their inputs and validity constraints. The
public PDF-malware study narrows this conceptual difference through constrained count
features, but supplies no Android result. Its experiments and the partner defence algorithms
retain their attribution. D7.4 is a confidential deliverable; the figures and empirical
comparisons in this chapter use public reports.

A focused subsequent study could compare an ordinary classifier with a training-selected
constant baseline and the same classifier under output thresholding. For detectors, paired
base predictions and alarms would test Equation~\ref{eq:advml:reversal-gain}. Budget curves
and per-example attack aggregation would then measure how the conclusions change under
validated perturbations. Those experiments remain separate from the historical evidence
presented here.

\section{Individual Contribution and Project Attribution}
\label{sec:advml:10}

The candidate's contribution is stated at section level, following his account of the
collaborative project. He wrote D7.1 Sections~2.2.2 and~2.3, and D7.2 Sections~2.1 and~2.2.
For D7.2 Sections~2.4 and~2.5 he contributed the concept and definitions. He edited and
refined partner text in D7.3, with no authored content claimed for that plan. He wrote
Sections~3.1 and~3.2 of D7.4, Chapter~5 of D7.5, and the whole of D7.6
\cite{sparta2021d71,sparta2020d72,sparta2021d73,sparta2021d75,sparta2022d76}.
His D7.4 contribution is recorded as part of this account, without using the confidential
report as an empirical source.

Within SAFAIR, the candidate contributed to organising the competition and to developing the
reusable benchmark framework. The contest toolkit and benchmark source code are published on
GitHub under the candidate's name
\cite{norouzian2021safairtoolkit,norouzian2021adversarialbenchmark}. The evaluated defensive
algorithms remain partner methods. In particular, the histopathology configurations originate
in the Vicomtech study of D7.2 Section~2.3.2; the PDF-malware curves belong to the partner
study in D7.5 Section~2.8. Authorship of the reporting deliverables does not transfer those
experiments or algorithms to the candidate. The mathematical analysis and source comparison
in this chapter examine their implications for the evaluation framework.

\section{Chapter Summary}
\label{sec:advml:11}

RQ4 concerns the translation of an adversary model into a reproducible benchmark and the
relation between clean performance and adversarial robustness. The SAFAIR work delivered a
competition with hidden evaluation samples and a reusable framework for evaluating trained
models against reference attacks. The framework established common component interfaces and
an execution procedure with adaptations for different prediction tasks. Its reported use by
Vicomtech on histopathology images documents an application separate from the competition.
The subsequent analysis identifies the experimental records and measurement checks required
to interpret the resulting scores.

The evidence supports several bounded conclusions. The histopathology table gives different
comparisons under accuracy loss and absolute attacked accuracy. The SD development account
reports improved transfer-attack performance in a configuration with lower clean accuracy,
but changes in ensemble size and training prevent causal attribution. The constant-prediction
result proves that the historical delta score can assign zero loss to a classifier with low
utility. Logit thresholding provides a separate mechanism: the attribute decisions remain
unchanged while direct gradients are obstructed. The invariant activation-detector accuracy
does not identify either mechanism.

These findings support an evaluation that reports clean and attacked utility together, tests
attacks against the full defence, and records individual predictions with valid perturbation
budgets. For detector-based systems, the response to an alarm is part of the method being
evaluated. The constrained PDF example further shows why the validity of an input must be
specified for its domain. The chapter establishes these methodological consequences from
the documented systems and their measurements, while leaving a general clean--robustness
relation and experiment-level reproduction unresolved.
